% !TeX encoding = UTF-8
%% \textbf{重庆大学}通用毕业论文\LaTeXe{}模板
%%% 使用前请先阅读使用文档和用户协议，内有详细介绍。Happy Texing! :)
%% =======================================================
\documentclass%
	[type=doctor, bilinguallist=apart,]{cquthesis}%
% 可用选项：
% type=[bachelor|master|doctor],      % 必选，毕业论文类型，以下项目不填时为默认
% liberalformat,                      % 可选，仅适用本科生，使用文学类论文标题格式，默认未打开
% proffesionalmaster=[true|false],    % 可选，仅适用研究生，是(true)否(false)专业硕士，默认为否
% printmode=[oneside|twoside|auto],	  % 可选，论文打印方式，默认采用auto按页数要求自动判定
% openany,|openright,                 % 可选，双面打印时每章的第一页仅右页开启，默认右页开启(openright)
% bilinguallist=[off|combined|apart], % 可选，图录表录等分别按双语题注混编(combined)，分开编录(apart)，默认关(off)
% blindtrail,                         % 可选，盲审模式，开启后封面姓名和致谢部分会隐藏，详情请参阅用户文档，默认关
% draft,                              % 写作期间可选，不渲染图片，关闭外围功能，加快预览速度，默认未开启

% 请在cquthesis.sty文件中定义其他会用到的宏包和自己的变量
% 这样可以防止main.tex太过臃肿。
\usepackage{cquthesis}

% 定义所有的图片文件在 figures 子目录下
\graphicspath{{figures/}}

%*** 写作时，使用这个命令只渲染你想查看的部分，提升工作效率，定稿时注释掉整行
%\includeonly{contents/experiment,contents/analysis,}

\begin{document}

\cqusetup{
%	************	注意	************
%	* 1. \cqusetup{}中不能出现全空的行，如果需要全空行请在行首注释
%	* 2. 不需要的配置信息可以放心地坐视不理、留空、删除或注释（都不会有影响）
%	*
%	********************************
% ===================
%	论文的中英文题目
% ===================
  ctitle = {\ce{C60}表面的纳米形貌的可控性转变\\兼论\LaTeX{}在论文排版中的运用},
  etitle = {To Use \LaTeX{} in the Typeseting of\\Graduating Work for CQU},
% ===================
% 作者部分的信息
% \secretize{}为盲审标记点，在打开盲审开关时内容会自动被替换为***输出，盲审开关默认关闭
% ===================
  cauthor = \secretize{李振楠},	% 你的姓名，以下每项都以英文逗号结束
  eauthor = \secretize{Zhennan~Li},	% 姓名拼音，~代表不会断行的空格
  studentid = \secretize{20128888},	% 仅本科生，学号
  csupervisor = \secretize{孙麟~~教授},	% 导师的姓名
  esupervisor = \secretize{{Prof.~Lin Sun}},	% 导师的姓名拼音
  cassistsupervisor = \secretize{}, % 本科生可选，助理指导教师姓名，不用时请留空为{}
  cextrasupervisor = \secretize{}, % 本科生可选，校外指导教师姓名，不用时请留空为{}
  eassistsupervisor = \secretize{}, % 本科生可选，助理指导教师或/和校外指导教师姓名拼音，不用时请留空为{}
  cpsupervisor = \secretize{丁小明~~工程师}, % 仅专硕，兼职导师姓名
  epsupervisor = \secretize{Eng.~Xiaoming~Ding},	% 仅专硕，兼职导师姓名拼音
  cclass = 工学,	% 博士生和学硕填学科门类，学硕填学科类型
  edgree = {Degree of Master of Enginnering},	% 专硕填Professional Degree，其他按实情填写
% 提示：如果内容太长，可以用\zihao{}命令控制字号，作用范围：{}内
  cmajor = {\zihao{-4}材料科学与工程（材料科学专业方向）},	% 专硕不需填，填写专业名称
  emajor = Material Science and Engineering, % % 专硕不需填，填写专业英文名称
% ===================
% 底部的学院名称和日期
% ===================
  cdepartment = 材料科学与工程学院,	%学院名称
  edepartment = College of Material Science and Engineering,	%学院英文名称
% ===================
% 封面的日期可以自动生成（注释掉时），也可以解除注释手动指定，例如：二〇一六年五月
% ===================
%	mycdate = {中文日期},
%	myedate = {Date in English},
}% End of \cqusetup
% ===================
%
% 论文的摘要
%
% ===================
\begin{cabstract}	% 中文摘要
	本文档是\cquthesis{}的示例文档。
		
	\cquthesis{}是重庆大学毕业论文的\LaTeX{}模板，支持学士、硕士、博士论文的排版。合理使用本模板可以大大减轻重庆大学毕业生在毕业论文撰写过程中的排版工作量。
	
	\cquthesis{}根据重庆大学《重庆大学本科设计（论文）撰写规范化要求（2007年修订版）》和《重庆大学博士、硕士论文撰写格式标准（2007年修订版）》编写，力求合规，简洁，易于实现，用户友好。
	
	本模板的特色：
  \begin{itemize}
  	\item 支持重庆大学本科（文学、理工）、硕士（学术、专业）、博士的毕业论文格式；
  	\item 内置封面、目录、索引、授权书等论文部件，可按需自动生成；
  	\item 自动侦测文档页数，生成相应的单面打印/双面打印PDF文件；
  	\item 预置一批优化过的宏包和小功能，包含国际标准单位、化学式支持、三线表等，可按需开启。
  \end{itemize}
\end{cabstract}
% 中文关键词，请使用英文逗号分隔：
\ckeywords{重庆大学,\LaTeX,\LaTeXe,论文,模板}

\begin{eabstract}	% 英文摘要
	LaTeX is a document preparation system for high-quality typesetting. It is most often used for medium-to-large technical or scientific documents but it can be used for almost any form of publishing.
	
  LaTeX contains features for:
\begin{enumerate}
  	\item Typesetting journal articles, technical reports, books, and slide presentations.
  	\item Control over large documents containing sectioning, cross-references, tables and figures.
  	\item Typesetting of complex mathematical formulas.
  	\item Advanced typesetting of mathematics with AMS-LaTeX.
  	\item Automatic generation of bibliographies and indexes.
  	\item Multi-lingual typesetting.
  	\item Inclusion of artwork, and process or spot colour.
  	\item Using PostScript or Metafont fonts.
  \end{enumerate}
  (Quote from \textit{https://latex-project.org/intro.html})  
\end{eabstract}
% 英文关键词，请使用英文逗号分隔，关键词内可以空格：
\ekeywords{bachelor, master, doctor, all support, white space is okay here}

% 封面和摘要配置完成

\makecover %%% 封面部分


\frontmatter %%%前置部分(封面后绪论前)
%% 摘要
\makeabstract
%% 目录，注意需要多次编译才能更新
\tableofcontents
%% 插图索引，可选，如不用可注释掉
\listoffigures
\listoffiguresEN
%% 表格索引，可选
\listoftables
\listoftablesEN
%% 公式索引，可选
\listofequations
\listofequationsEN
%% 符号对照表，可选
% !TeX root = ../main.tex

\begin{denotation}[3cm]
\item[HPC] 高性能计算 (High Performance Computing)
\item[cluster] 集群
\item[Itanium] 安腾
\item[SMP] 对称多处理
\item[API] 应用程序编程接口
\item[PI] 聚酰亚胺
\item[MPI] 聚酰亚胺模型化合物，N-苯基邻苯酰亚胺
\item[PBI] 聚苯并咪唑
\item[MPBI] 聚苯并咪唑模型化合物，N-苯基苯并咪唑
\item[PY] 聚吡咙
\item[PMDA-BDA]	均苯四酸二酐与联苯四胺合成的聚吡咙薄膜
\item[$\increment G$] 活化自由能 (Activation Free Energy)
\item[$\chi$] 传输系数 (Transmission Coefficient)
\item[$E$] 能量
\item[$m$] 质量
\item[$c$] 光速
\item[$P$] 概率
\item[$T$] 时间
\item[$v$] 速度
\item[劝学] 君子曰：学不可以已。青，取之于蓝，而青于蓝；冰，水为之，而寒于水。木
  直中绳。輮以为轮，其曲中规。虽有槁暴，不复挺者，輮使之然也。故木受绳则直，金就
  砺则利，君子博学而日参省乎己，则知明而行无过矣。吾尝终日而思矣，不如须臾之所学
  也；吾尝跂而望矣，不如登高之博见也。登高而招，臂非加长也，而见者远；顺风而呼，
  声非加疾也，而闻者彰。假舆马者，非利足也，而致千里；假舟楫者，非能水也，而绝江
  河，君子生非异也，善假于物也。积土成山，风雨兴焉；积水成渊，蛟龙生焉；积善成德，
  而神明自得，圣心备焉。故不积跬步，无以至千里；不积小流，无以成江海。骐骥一跃，
  不能十步；驽马十驾，功在不舍。锲而舍之，朽木不折；锲而不舍，金石可镂。蚓无爪牙
  之利，筋骨之强，上食埃土，下饮黄泉，用心一也。蟹六跪而二螯，非蛇鳝之穴无可寄托
  者，用心躁也。—— 荀况
\end{denotation}



% % 也可以使用 nomencl 宏包：

% \printnomenclature[3cm]

% \nomenclature{HPC}{高性能计算 (High Performance Computing)}
% \nomenclature{cluster}{集群}
% \nomenclature{Itanium}{安腾}
% \nomenclature{SMP}{对称多处理}
% \nomenclature{API}{应用程序编程接口}
% \nomenclature{PI}{聚酰亚胺}
% \nomenclature{MPI}{聚酰亚胺模型化合物，N-苯基邻苯酰亚胺}
% \nomenclature{PBI}{聚苯并咪唑}
% \nomenclature{MPBI}{聚苯并咪唑模型化合物，N-苯基苯并咪唑}
% \nomenclature{PY}{聚吡咙}
% \nomenclature{PMDA-BDA}{均苯四酸二酐与联苯四胺合成的聚吡咙薄膜}
% \nomenclature{$\increment G$}{活化自由能 (Activation Free Energy)}
% \nomenclature{$\chi$}{传输系数 (Transmission Coefficient)}
% \nomenclature{$E$}{能量}
% \nomenclature{$m$}{质量}
% \nomenclature{$c$}{光速}
% \nomenclature{$P$}{概率}
% \nomenclature{$T$}{时间}
% \nomenclature{$v$}{速度}



\mainmatter %%% 主体部分(绪论开始，结论为止)
%* 子文件的多少和内容由你决定（最好以章为单位），基本原则是提速预览、脉络清晰、管理容易。


\documentclass[12pt]{report}
\usepackage[utf8]{inputenc}
\usepackage[brazilian,brazil]{babel}
\usepackage{fancyhdr,setspace,float,graphicx,lscape,array,longtable,colortbl,amsmath,amssymb,booktabs,multirow,hyperref,pdfpages,tocloft,titlesec,lipsum,natbib}
\usepackage[sectionbib]{chapterbib}
\begin{document}
%***
\clearpage
\linenumbers % numeração de linhas
\modulolinenumbers[3] % numeração de linhas
%***
% Texto
%***
\chapter{Nam dui ligula}
%***
\lipsum[2-2] Nam dui ligula, fringilla a, euismod sodales, sollicitudin vel, wisi. Morbi
auctor lorem non justo \citep{lamport1986latex}.
\section{Euismod sodales}
\lipsum[2-3]
%*** APAGAR O EXEMPLO ACIMA







%*** REFERÊNCIAS
\bibliography{referencias.bib}
\bibliographystyle{apalike}
\end{document}

\chapter{实验参数和流程}
\section{5分钟语法参考}

{\kaishu 要流畅使用\cquthesis 需要用户对\LaTeX 以及\textsc{Bib}\TeX 有一定了解，下面这个语法参考只能起到抛砖引玉的作用。如果你从来没有接触过\LaTeX 或者\XeTeX ，建议先学习相关知识，磨刀不误砍柴工。}
\medskip
\begin{itemize}
	\item \LaTeX 源文件中，主要有三种元素：你的文字，命令，以及环境；
	\item 直接输入即可你想要写的文字即可，对于英文，文字间多于一个的空格都会转为一个空格；
	\item 如果你想开启一个新的自然段，请在写新内容前空一个（或多个）全空的行；
	\item \LaTeX 的命令{\heiti 全部}都以\cs{ }开头，例如\cs{XeTeX}可以生成\XeTeX ；
	\item 有的命令{\heiti 必须}带参数，比如\cs{zihao\{-4\}}可以将命令之后的内容的字号调整为小四；
	\item 有的命令能带可选参数，例如\cs{usingpackage\{metalogo\}}可以载入\pkg{metalogo}宏包；
	\item 宏包中有宏包作者自己定义的命令，能够让你更容易地完成某些事情，比如\pkg{mhchem}能够引入让你方便地表示化学式的命令\cs{ce}；
	\item \LaTeX 的源代码主要分为两个部分，导言部分和文档部分。其中，文档部分以\cs{begin\{document\}}开头，以\cs{end\{document\}}结尾，只有在这个范围内你才能完成排版工作；
	\item \LaTeX 对（简单或复杂的）数学式的支持是其一大亮点，数学环境使用\texttt{\${ }\$}包裹；
	\item 环境由\csgo{begin}{环境名}开头，以\csgo{end}{环境名}结尾，是的，文档部分是一个巨大的环境；
	\item 报错说没有这个命令？检查是否载入了必要的宏包，再检查命令后面是否直接跟随了汉字，在它们之间加个空格就好；
	\item \LaTeX 是一门语言，新手经常会遇到无法编译通过的语法错误，这时建议你仔细检查花括号是否平衡，命令是否敲错，参数数目和类型是否正确，如果还是不行，可以在网络上搜索一番或者问问同事。
	\item 命令之间或者之内的空格和缩进以及回车不是必须的，事实上没有它们\LaTeX 也可以正常工作，但是代码的可读性就会大打折扣了；
	\item 对了，使用\texttt{\%}来开启一个行注释，注释的内容不参与编译，你可以在这里写下自己的小秘密；
	\item 有质量的国内\TeX 社区是\textsc{CTeX}社区，更有质量的国外的是\textsf{StackExchange}；
	\item \TeX\textsc{Studio}是一个很棒的\LaTeX 编辑环境，推荐你尝试一番。
\end{itemize}

\section{查询文档}
在你对宏包或者环境包有疑问的时候，可以再命令行中输入：

\texttt{texdoc 宏包或环境名称}

回车后对应的用户文档会自动打开。






\chapter{实验结果和分析}

\section{字体命令}\label{txt:FreqCmd}
{\kaishu 玲珑骰子安红豆，入骨相思知不知。\hfill ——温庭筠}
	
{\fangsong 愿得一心人，白头不相离。\hfill ——卓文君}
		
{\ifcsname youyuan\endcsname\youyuan\else[无 \cs{youyuan} 字体。]\fi 去年今日此门中，人面桃花相映红。\hfill ——崔护}
			
{\heiti 入我相思门，知我相思苦。\hfill ——李白}
				
{\ifcsname lishu\endcsname\lishu\else[无 \cs{lishu} 字体。]\fi 此情可待成追忆？只是当时已惘然。\hfill ——李商隐}
					
{\songti 雨打梨花深闭门，忘了青春，误了青春。\hfill ——唐寅}

使用\cs{textbf}和\cs{textit}以及\cs{underline}的效果分别如下：

这句话的\textbf{文字}分别\textit{使用}了三种命令来\underline{处理}。

The \textbf{words} in this sentences are \textit{processed} with three different \underline{cmd}.

\section{表格样本}

\subsection{基本表格}
\label{sec:basictable}

模板中关于表格的宏包有三个： \pkg{booktabs}、\pkg{array} 和\pkg{longtabular}。三线表可以用 \pkg{booktabs}提供的 \cs{toprule}、\cs{midrule} 和 \cs{bottomrule}。它们与\pkg{longtable} 能很好的配合使用。
\begin{table}[htb]
	\centering
	\begin{minipage}[t]{0.9\linewidth} % 如果想在表格中使用脚注，minipage是个不错的办法
	\caption[模板文件]{模板文件。如果表格的标题很长，那么在表格索引中就会很不美观，所以要像 chapter 那样在前面用中括号写一个简短的标题。这个标题会出现在索引中。}
	\label{tab:template-files}
	\begin{tabularx}{\linewidth}{lX}
		\toprule
		{\heiti 文件名} & {\heiti 描述} \\
		\midrule
		cquthesis.cls & 模板类文件\footnote{这是一个脚注}\\
		cquthesis.cfg & 模板配置文件\footnote{这是又一个脚注}\\
		cqunumberical.bst & 参考文献 BIB\TeX\ 样式文件。\\
		cquthesis.sty & 常用的包和命令写在这里，减轻主文件的负担。\footnote{同一页上的脚注最多支持到10个}\\
		\bottomrule
		\end{tabularx}
	\end{minipage}
\end{table}

首先来看一个最简单的表格。\autoref{tab:template-files} 列举了本模板主要文件及其功能。请大家注意三线表中各条线对应的命令。这个例子还展示了如何在表格中正确使用脚注。由于 \LaTeX{} 本身不支持在表格中使用\cs{footnote}，所以我们不得不将表格放在小页中，而且最好将表格的宽度设置为小页的宽度，这样脚注看起来才更美观。

\subsection{双语题注和复杂表格}
\label{sec:complicatedtable}
使用\cs{bicaption}\marg{中文}\marg{英文}可以对图或者表的浮动体添加双语题注，对方程式进行双语题注，请使用\cs{eqlist}\marg{中文}\oarg{英文}，注意括号。

我们经常会在表格下方标注数据来源，或者对表格里面的条目进行解释。前面的脚注是一种不错的方法，如果不喜欢脚注，可以在表格后面写注释，比如\autoref{tab:tabexamp1}。
\begin{table}[htbp]
	\centering
	\bicaption{复杂表格示例}{A more structured table}
	\label{tab:tabexamp1}
	\begin{minipage}[t]{0.8\textwidth} 
	\begin{tabularx}{\linewidth}{|l|X|X|X|X|}
		\hline
		\multirow{2}*{\diagbox[width=5em]{x}{y}} & \multicolumn{2}{c|}{First Half} & \multicolumn{2}{c|}{Second Half}\\\cline{2-5}
		& 1st Qtr &2nd Qtr&3rd Qtr&4th Qtr \\ \hline
		East$^{*}$ &   20.4&   27.4&   90&     20.4 \\
		West$^{**}$ &   30.6 &   38.6 &   34.6 &  31.6 \\ \hline
	\end{tabularx}\\[2pt]
	\footnotesize 
	*：东部\\
	**：西部
	\end{minipage}
\end{table}

此外，表~\ref{tab:tabexamp1} 同时还演示了另外两个功能：1）通过 \pkg{tabularx} 的\texttt{|X|} 扩展实现表格自动放大；2）通过命令 \cs{diagbox} 在表头部分插入反斜线。

\begin{table}[htbp]
	\noindent\begin{minipage}{0.5\textwidth}
		\centering
		\caption{第一个并排子表格}
		\label{tab:parallel1}
		\begin{tabular}{p{2cm}p{2cm}}
					\toprule
					No. & Name \\\midrule
					\xuhao[1] & Fox \\
					\xuhao & Panda \\
					\xuhao & Dog \\
					\bottomrule
		\end{tabular}
	\end{minipage}%
	\setxuhao[2]
	\begin{minipage}{0.5\textwidth}
		\centering
		\bicaption{第二个并排子表格}{The second subtable in one row}
		\label{tab:parallel2}
		\begin{tabular}{p{2cm}p{2cm}}
			\toprule
			No. & Name \\\midrule
			\xuhao[1] & Charlie \\
			\xuhao & Jack \\
			\xuhao & Tom \\
			\bottomrule
		\end{tabular}
	\end{minipage}
\end{table}

\begin{table}[htbp]
	\centering
	\caption{并排子表格}
	\label{tab:subtable}
	\subcaptionbox{第一个子表格}
	{
		\begin{tabular}{p{2cm}p{2cm}}
			\toprule
			111 & 222 \\\midrule
			222 & 333 \\\bottomrule
		\end{tabular}
	}
	\hskip2cm
	\subcaptionbox{第二个子表格}
	{
		\begin{tabular}{p{2cm}p{2cm}}
			\toprule
			111 & 222 \\\midrule
			222 & 333 \\\bottomrule
		\end{tabular}
	}
\end{table}

不可否认 \LaTeX{} 的表格功能没有想象中的那么强大，不过只要足够认真，足够细致，同样可以排出来非常复杂非常漂亮的表格。

\tabref{tab:parallel1}和\tabref{tab:parallel2}展示了\cs{xuhao}和\cs{xuhao}\texttt{[1]}的使用，可以达到自动编号的效果。不过要记得在每次使用之前使用\cs{resetxuhao}，或者\cs{xuhao}\texttt{[1]}。使用\cs{setxuhao}\oarg{1-6}可以更改序号的标记方式，如\tabref{tab:parallel2}所示。详细用法请参阅用户手册。

\begin{longtable}[c]{c*{6}{r}}
	\bicaption[实验数据]{实验数据，这个题注是双语的，而且十分的长，注意这在索引中的处理方式}[Data in experiment]{Data in experiment, and this is a really long long long long long long long long long long text.}\label{tab:performance}\\
	\toprule
	测试程序 & \multicolumn{1}{c}{正常运行} & \multicolumn{1}{c}{同步} & \multicolumn{1}{c}{检查点} & \multicolumn{1}{c}{卷回恢复}
	& \multicolumn{1}{c}{进程迁移} & \multicolumn{1}{c}{检查点} \\
	& \multicolumn{1}{c}{时间 (s)}& \multicolumn{1}{c}{时间 (s)}&
	\multicolumn{1}{c}{时间 (s)}& \multicolumn{1}{c}{时间 (s)}& \multicolumn{1}{c}{
		时间 (s)}&  文件（KB）\\\midrule
	\endfirsthead
	\multicolumn{7}{c}{续表~\thetable\hskip1em 实验数据}\\
	\toprule
	测试程序 & \multicolumn{1}{c}{正常运行} & \multicolumn{1}{c}{同步} & \multicolumn{1}{c}{检查点} & \multicolumn{1}{c}{卷回恢复}
	& \multicolumn{1}{c}{进程迁移} & \multicolumn{1}{c}{检查点} \\
	& \multicolumn{1}{c}{时间 (s)}& \multicolumn{1}{c}{时间 (s)}&
	\multicolumn{1}{c}{时间 (s)}& \multicolumn{1}{c}{时间 (s)}& \multicolumn{1}{c}{
		时间 (s)}&  文件（KB）\\\midrule
	\endhead
	\hline
	\multicolumn{7}{r}{续下页}
	\endfoot
	\endlastfoot
	CG.A.2 & 23.05 & 0.002 & 0.116 & 0.035 & 0.589 & 32491 \\
	CG.A.4 & 15.06 & 0.003 & 0.067 & 0.021 & 0.351 & 18211 \\
	CG.A.8 & 13.38 & 0.004 & 0.072 & 0.023 & 0.210 & 9890 \\
	CG.B.2 & 867.45 & 0.002 & 0.864 & 0.232 & 3.256 & 228562 \\
	CG.B.4 & 501.61 & 0.003 & 0.438 & 0.136 & 2.075 & 123862 \\
	CG.B.8 & 384.65 & 0.004 & 0.457 & 0.108 & 1.235 & 63777 \\
	MG.A.2 & 112.27 & 0.002 & 0.846 & 0.237 & 3.930 & 236473 \\
	MG.A.4 & 59.84 & 0.003 & 0.442 & 0.128 & 2.070 & 123875 \\
	MG.A.8 & 31.38 & 0.003 & 0.476 & 0.114 & 1.041 & 60627 \\
	MG.B.2 & 526.28 & 0.002 & 0.821 & 0.238 & 4.176 & 236635 \\
	MG.B.4 & 280.11 & 0.003 & 0.432 & 0.130 & 1.706 & 123793 \\
	MG.B.8 & 148.29 & 0.003 & 0.442 & 0.116 & 0.893 & 60600 \\
	LU.A.2 & 2116.54 & 0.002 & 0.110 & 0.030 & 0.532 & 28754 \\
	LU.A.4 & 1102.50 & 0.002 & 0.069 & 0.017 & 0.255 & 14915 \\
	LU.A.8 & 574.47 & 0.003 & 0.067 & 0.016 & 0.192 & 8655 \\
	LU.B.2 & 9712.87 & 0.002 & 0.357 & 0.104 & 1.734 & 101975 \\
	LU.B.4 & 4757.80 & 0.003 & 0.190 & 0.056 & 0.808 & 53522 \\
	LU.B.8 & 2444.05 & 0.004 & 0.222 & 0.057 & 0.548 & 30134 \\
	CG.B.2 & 867.45 & 0.002 & 0.864 & 0.232 & 3.256 & 228562 \\
	CG.B.4 & 501.61 & 0.003 & 0.438 & 0.136 & 2.075 & 123862 \\
	CG.B.8 & 384.65 & 0.004 & 0.457 & 0.108 & 1.235 & 63777 \\
	MG.A.2 & 112.27 & 0.002 & 0.846 & 0.237 & 3.930 & 236473 \\
	MG.A.4 & 59.84 & 0.003 & 0.442 & 0.128 & 2.070 & 123875 \\
	MG.A.8 & 31.38 & 0.003 & 0.476 & 0.114 & 1.041 & 60627 \\
	MG.B.2 & 526.28 & 0.002 & 0.821 & 0.238 & 4.176 & 236635 \\
	MG.B.4 & 280.11 & 0.003 & 0.432 & 0.130 & 1.706 & 123793 \\
	MG.B.8 & 148.29 & 0.003 & 0.442 & 0.116 & 0.893 & 60600 \\
	LU.A.2 & 2116.54 & 0.002 & 0.110 & 0.030 & 0.532 & 28754 \\
	LU.A.4 & 1102.50 & 0.002 & 0.069 & 0.017 & 0.255 & 14915 \\
	LU.A.8 & 574.47 & 0.003 & 0.067 & 0.016 & 0.192 & 8655 \\
	LU.B.2 & 9712.87 & 0.002 & 0.357 & 0.104 & 1.734 & 101975 \\
	LU.B.4 & 4757.80 & 0.003 & 0.190 & 0.056 & 0.808 & 53522 \\
	LU.B.8 & 2444.05 & 0.004 & 0.222 & 0.057 & 0.548 & 30134 \\
	EP.A.2 & 123.81 & 0.002 & 0.010 & 0.003 & 0.074 & 1834 \\
	EP.A.4 & 61.92 & 0.003 & 0.011 & 0.004 & 0.073 & 1743 \\
	EP.A.8 & 31.06 & 0.004 & 0.017 & 0.005 & 0.073 & 1661 \\
	EP.B.2 & 495.49 & 0.001 & 0.009 & 0.003 & 0.196 & 2011 \\
	EP.B.4 & 247.69 & 0.002 & 0.012 & 0.004 & 0.122 & 1663 \\
	EP.B.8 & 126.74 & 0.003 & 0.017 & 0.005 & 0.083 & 1656 \\
	\bottomrule
\end{longtable}

如果你要排版的表格长度超过一页，那么推荐使用 \pkg{longtable} 或者 \pkg{supertabular}宏包，模板对 \pkg{longtable} 进行了相应的设置，所以用起来可能简单一些。表~\ref{tab:performance} 就是 \pkg{longtable} 的简单示例。

\section{定理环境}
\label{sec:theorem}

给大家演示一下各种和证明有关的环境：

\begin{assumption}
	假设以下数学方程成立：
	\begin{eqnarray}
	\label{eq:eqnxmp}
	c & = & a^2 - b^2\\
	& = & (a+b)(a-b)
	\end{eqnarray}
\end{assumption}

\begin{assumption}
	依然假设以下数学方程成立，注意整个系统是自动编号的：
	\begin{eqnarray}
	\label{eq:eqnxmp2}
	c & = & a^2 - b^2\\
	& = & (a+b)(a-b)
	\end{eqnarray}
\end{assumption}

\begin{definition}
	我们定义\ref{eq:eqnxmp}中的方程名称为\cquthesis 。你看，环境里是可以相互引用的。
\end{definition}

\begin{proposition}
	曾子曰：「吾日三省吾身 —— 为人谋而不忠乎？与朋友交而不信乎？传不习乎？」
\end{proposition}

多么凄美的命题啊！其日牛马嘶，新妇入青庐，奄奄黄昏后，寂寂人定初，我命绝今日，
魂去尸长留，揽裙脱丝履，举身赴清池，府吏闻此事，心知长别离，徘徊庭树下，自挂东南
枝。

\begin{remark}
	天不言自高，水不言自流。
	\begin{gather*}
	\begin{split} 
	\varphi(x,z)
	&=z-\gamma_{10}x-\gamma_{mn}x^mz^n\\
	&=z-Mr^{-1}x-Mr^{-(m+n)}x^mz^n
	\end{split}\\[6pt]
	\begin{align} \zeta^0&=(\xi^0)^2,\\
	\zeta^1 &=\xi^0\xi^1,\\
	\zeta^2 &=(\xi^1)^2,
	\end{align}
	\end{gather*}
\end{remark}

天尊地卑，乾坤定矣。卑高以陈，贵贱位矣。 动静有常，刚柔断矣。方以类聚，物以群分，
吉凶生矣。在天成象，在地成形，变化见矣。鼓之以雷霆，润之以风雨，日月运行，一寒一
暑，乾道成男，坤道成女。乾知大始，坤作成物。乾以易知，坤以简能。易则易知，简则易
从。易知则有亲，易从则有功。有亲则可久，有功则可大。可久则贤人之德，可大则贤人之
业。易简，而天下矣之理矣；天下之理得，而成位乎其中矣。

\begin{axiom}
	两点间直线段距离最短。  
	\begin{align}
	x&\equiv y+1\pmod{m^2}\\
	x&\equiv y+1\mod{m^2}\\
	x&\equiv y+1\pod{m^2}
	\end{align}
\end{axiom}

《彖曰》：大哉乾元，万物资始，乃统天。云行雨施，品物流形。大明始终，六位时成，时
乘六龙以御天。乾道变化，各正性命，保合大和，乃利贞。首出庶物，万国咸宁。

《象曰》：天行健，君子以自强不息。潜龙勿用，阳在下也。见龙再田，德施普也。终日乾
乾，反复道也。或跃在渊，进无咎也。飞龙在天，大人造也。亢龙有悔，盈不可久也。用九，
天德不可为首也。 　　

\begin{lemma}
	《猫和老鼠》是我最爱看的动画片。
	\begin{multline*}%\tag*{[a]} % 这个不出现在索引中
	\int_a^b\biggl\{\int_a^b[f(x)^2g(y)^2+f(y)^2g(x)^2]
	-2f(x)g(x)f(y)g(y)\,dx\biggr\}\,dy \\
	=\int_a^b\biggl\{g(y)^2\int_a^bf^2+f(y)^2
	\int_a^b g^2-2f(y)g(y)\int_a^b fg\biggr\}\,dy
	\end{multline*}
\end{lemma}

行行重行行，与君生别离。相去万余里，各在天一涯。道路阻且长，会面安可知。胡马依北
风，越鸟巢南枝。相去日已远，衣带日已缓。浮云蔽白日，游子不顾返。思君令人老，岁月
忽已晚。  弃捐勿复道，努力加餐饭。

\begin{theorem}\label{the:theorem1}
	犯我强汉者，虽远必诛\hfill —— 陈汤（汉）
\end{theorem}
\begin{subequations}
	\begin{align}
	y & = 1 \\
	y & = 0
	\end{align}
\end{subequations}
道可道，非常道。名可名，非常名。无名天地之始；有名万物之母。故常无，欲以观其妙；
常有，欲以观其徼。此两者，同出而异名，同谓之玄。玄之又玄，众妙之门。上善若水。水
善利万物而不争，处众人之所恶，故几于道。曲则全，枉则直，洼则盈，敝则新，少则多，
多则惑。人法地，地法天，天法道，道法自然。知人者智，自知者明。胜人者有力，自胜
者强。知足者富。强行者有志。不失其所者久。死而不亡者寿。

\begin{proof}
	燕赵古称多感慨悲歌之士。董生举进士，连不得志于有司，怀抱利器，郁郁适兹土，吾
	知其必有合也。董生勉乎哉？
	
	夫以子之不遇时，苟慕义强仁者，皆爱惜焉，矧燕、赵之士出乎其性者哉！然吾尝闻
	风俗与化移易，吾恶知其今不异于古所云邪？聊以吾子之行卜之也。董生勉乎哉？
	
	吾因子有所感矣。为我吊望诸君之墓，而观于其市，复有昔时屠狗者乎？为我谢
	曰：“明天子在上，可以出而仕矣！” \hfill —— 韩愈《送董邵南序》
\end{proof}

\begin{corollary}
	四川话配音的《猫和老鼠》是世界上最好看最好听最有趣的动画片。
	\begin{alignat}{3}
	V_i & =v_i - q_i v_j, & \qquad X_i & = x_i - q_i x_j,
	& \qquad U_i & = u_i,
	\qquad \text{for $i\ne j$;}\label{eq:B}\\
	V_j & = v_j, & \qquad X_j & = x_j,
	& \qquad U_j & u_j + \sum_{i\ne j} q_i u_i.
	\end{alignat}
\end{corollary}

迢迢牵牛星，皎皎河汉女。
纤纤擢素手，札札弄机杼。
终日不成章，泣涕零如雨。
河汉清且浅，相去复几许。
盈盈一水间，脉脉不得语。

\begin{example}
	大家来看这个例子。
	\begin{equation}
	\label{ktc}
	\left\{\begin{array}{l}
	\nabla f({\mbox{\boldmath $x$}}^*)-\sum\limits_{j=1}^p\lambda_j\nabla g_j({\mbox{\boldmath $x$}}^*)=0\\[0.3cm]
	\lambda_jg_j({\mbox{\boldmath $x$}}^*)=0,\quad j=1,2,\cdots,p\\[0.2cm]
	\lambda_j\ge 0,\quad j=1,2,\cdots,p.
	\end{array}\right.
	\end{equation}
\end{example}

\begin{exercise}
	清列出 Andrew S. Tanenbaum 和 W. Richard Stevens 的所有著作。
\end{exercise}

\begin{conjecture} \textit{Poincare Conjecture} If in a closed three-dimensional
	space, any closed curves can shrink to a point continuously, this space can be
	deformed to a sphere.
\end{conjecture}

\begin{problem}
	回答还是不回答，是个问题。 
\end{problem}

如何引用定理~\ref{the:theorem1} 呢？加上 \cs{label} 使用 \cs{ref} 即可。

\section{参考文献}
\label{sec:bib}
重庆大学的要求是参考文献以上标的形式标注于论述之后，就像这样：

研究表明\cite{r1}，早睡早起有益身体健康。如果想同时引用多个文献\cite{r2,r3,r4,r6}，只需要在\csgo{cite}{\null}中用逗号分开\textsf{citeKey}就好。

\cquthesis 同时提供正文模式的参考文献引用功能\cs{inlinecite}，适用于以下情况：

文献\inlinecite{r6,z1,z2,z3}表明，文献\inlinecite{r7,r8,r9,r10}所述的情况是有理论依据的。

\section{数学公式}
\label{sec:equation}
贝叶斯公式如式~(\ref{equ:chap1:bayes})，其中 $p(y|\mathbf{x})$ 为后验；
$p(\mathbf{x})$ 为先验；分母 $p(\mathbf{x})$ 为归一化因子。
\begin{equation}
\label{equ:chap1:bayes}
p(y|\mathbf{x}) = \frac{p(\mathbf{x},y)}{p(\mathbf{x})}=
\frac{p(\mathbf{x}|y)p(y)}{p(\mathbf{x})} 
\end{equation}

论文里面公式越多，\TeX{} 就越 happy。再看一个 \pkg{amsmath} 的例子：
\newcommand{\envert}[1]{\left\lvert#1\right\rvert} 
\begin{equation}\label{detK2}
\det\mathbf{K}(t=1,t_1,\dots,t_n)=\sum_{I\in\mathbf{n}}(-1)^{\envert{I}}
\prod_{i\in I}t_i\prod_{j\in I}(D_j+\lambda_jt_j)\det\mathbf{A}
^{(\lambda)}(\overline{I}|\overline{I})=0.
\end{equation} 

前面定理示例部分列举了很多公式环境，可以说把常见的情况都覆盖了，大家在写公式的时候一定要好好看 \pkg{amsmath} 的文档，并参考模板中的用法：
\begin{multline*}%\tag{[b]} % 这个出现在索引中的
\int_a^b\biggl\{\int_a^b[f(x)^2g(y)^2+f(y)^2g(x)^2]
-2f(x)g(x)f(y)g(y)\,dx\biggr\}\,dy \\
=\int_a^b\biggl\{g(y)^2\int_a^bf^2+f(y)^2
\int_a^b g^2-2f(y)g(y)\int_a^b fg\biggr\}\,dy
\end{multline*}

这里还有一个多级规划公式，这个公式使用\csgo{listeq}{索引名}手动加入了目录后的索引。
\begin{equation}\label{bilevel}
\left\{\begin{array}{l}
\max\limits_{{\mbox{\footnotesize\boldmath $x$}}} F(x,y_1^*,y_2^*,\cdots,y_m^*)\\[0.2cm]
\mbox{subject to:}\\[0.1cm]
\qquad G(x)\le 0\\[0.1cm]
\qquad(y_1^*,y_2^*,\cdots,y_m^*)\mbox{ solves problems }(i=1,2,\cdots,m)\\[0.1cm]
\qquad\left\{\begin{array}{l}
\max\limits_{{\mbox{\footnotesize\boldmath $y_i$}}}f_i(x,y_1,y_2,\cdots,y_m)\\[0.2cm]
\mbox{subject to:}\\[0.1cm]
\qquad g_i(x,y_1,y_2,\cdots,y_m)\le 0.
\end{array}\right.
\end{array}\right.
\end{equation}\listeq{多级规划公式}
这些跟规划相关的公式都来自于清华大学刘宝碇老师《不确定规划》的课件。以上的许多例子由清华大学的薛瑞尼同学编写。

\section{化学方程式}

使用\pkg{mhchem}的\csgo{ce}{化学式或方程式}能够让你很容易地表示出各种化学式和化学方程：

例如：
\begin{center}
	\ce{C6H5-CHO}\\ \ce{A\bond{~--}B\bond{~=}C\bond{-~-}D}\\ \ce{SO4^2- + Ba^2+ -> BaSO4 v}
\end{center}

复杂一点的方程式也不在话下，如\eqref{eq:chem}：
\begin{equation}\label{eq:chem}
	\ce{Zn^2+
		<=>[+ 2OH-][+ 2H+]
		$\underset{\text{amphoteres Hydroxid}}{\ce{Zn(OH)2 v}}$ <=>[+ 2OH-][+ 2H+]
		$\underset{\text{Hydroxozikat}}{\ce{[Zn(OH)4]^2-}}$
	}
\end{equation}\eqlist{复杂的化学方程式}[A sophisticated chemical equation]

这个方程式嵌套在了\pkg{equation}环境中，可用\cs{eqlist}(\cs{listeq}的别名，作用相同)来编排到索引中。

如果你需要一次列举多个化学式，可以用\cs{cec}命令，例如，\csgo{cec}{H2O,HCl,CCl4}的输出为\cec{H2O,HCl,CCl4}。

\section{国际单位制(SI Unit)}

\cquthesis 采用\pkg{siunitx}作为国际单位制支持宏包，以下是一些使用例子，这个包的文档写得非常不错，请在命令行里输入\texttt{texdoc siunitx}察看。
\begin{center}
	\num{.3e45}\\
	\num{1.654 x 2.34 x 3.430}\\
	\si{\kilogram\metre\per\second}\\    
	\SIlist{0.13;0.67;0.80}{\milli\metre}
\end{center}


\section{绘图}
\label{sec:draw}

本模板不预先装载任何绘图包（如 \pkg{pstricks，pgf} 等），完全由用户来决定。个人觉得 \pkg{pgf} 不错，不依赖于 Postscript。此外还有很多针对 \LaTeX{} 的GUI 作图工具，如 XFig(jFig), WinFig, Tpx, Ipe, Dia, Inkscape, LaTeXPiX,jPicEdt, jaxdraw 等等。

\section{插图}
\label{sec:graphs}

推荐《\LaTeXe\ 插图指南》。关于子图形的使用细节请参看 \pkg{subcaption} 宏包的说明文档。

\subsection{一个图形}
\label{sec:onefig}
一般图形都是处在浮动环境中。之所以称为浮动是指最终排版效果图形的位置不一定与源文
件中的位置对应\footnote{这是\LaTeX 的一个设计特性。}，这也是刚使
用 \LaTeX{} 同学可能遇到的问题。如果要强制固定浮动图形的位置，请使用 \pkg{float} 宏包，
它提供了 \texttt{[H]} 参数，比如图~\ref{fig:xfig1}。
\begin{figure}[htb] % use float package if you want it here
	\centering
	\includegraphics[height=4cm]{CQUbadge.pdf}
	\bicaption{重庆大学校徽}{Chongqing University badage}
	\label{fig:xfig1}
\end{figure}

大学之道，在明明德，在亲民，在止于至善。知止而后有定；定而后能静；静而后能安；安
而后能虑；虑而后能得。物有本末，事有终始。知所先后，则近道矣。古之欲明明德于天
下者，先治其国；欲治其国者，先齐其家；欲齐其家者，先修其身；欲修其身者，先正其心；
欲正其心者，先诚其意；欲诚其意者，先致其知；致知在格物。物格而后知至；知至而后
意诚；意诚而后心正；心正而后身 修；身修而后家齐；家齐而后国治；国治而后天下
平。自天子以至于庶人，壹是皆以修身为本。其本乱而未治者 否矣。其所厚者薄，而其所
薄者厚，未之有也！

\hfill —— 《大学》


\subsection{多个图形}
\label{sec:multifig}

如果多个图形相互独立，并不共用一个图形计数器，那么用 \texttt{minipage} 或者\texttt{parbox} 就可以。否则，请参看
图~\ref{fig:big1-subcaptionbox}，它包含两个小图，分别是图~\ref{fig:subfig1}和图~\ref{fig:subfig2}。推荐使用\cs{subcaptionbox}，因为可以像图~\ref{fig:big1-subcaptionbox} 那样对齐子图的标题，也可以使用\pkg{subcaption}宏包的\cs{subcaption}（放在minipage中，用法同\cs{caption}）或是\pkg{subfigure}、\pkg{subtable}环境，像图~\ref{fig:big1-subfigure}，不要再用 \cs{subfloat}、\cs{subfigure} 和 \cs{subtable}。

\begin{figure}[h]
	\centering%
	\subcaptionbox{第一个小图形\label{fig:subfig1}}[3cm] %标题的长度，超过则会换行，如下一个小图。
	{\includegraphics[height=4cm]{CQUbadge.pdf}}%
	\hspace{4em}%
	\subcaptionbox{第二个小图形，注意这个图略矮些。如果标题很长的话，它会自动换行\label{fig:subfig2}}
	{\includegraphics[height=3cm]{CQUbadge.pdf}}
	\caption{包含子图形的大图形(subcaptionbox示例)}
	\label{fig:big1-subcaptionbox}
\end{figure}
\begin{figure}[ht]
	\centering%
	\begin{subfigure}{3cm}
		\includegraphics[height=4cm]{CQUbadge.pdf}
		\caption{第一个小图形}
	\end{subfigure}%
	\hspace{4em}%
	\begin{subfigure}{0.5\textwidth}
		\includegraphics[height=3cm]{CQUbadge.pdf}
		\caption{第二个小图形，注意这个图略矮些。subfigure中同一行的子图在顶端对齐。}
	\end{subfigure}
	\caption{包含子图形的大图形(subfigure示例)}
	\label{fig:big1-subfigure}
\end{figure}

如果要把编号的两个图形并排，那么小页就非常有用了。
\begin{figure}
	\begin{minipage}{0.48\textwidth}
		\centering
		\includegraphics[height=5cm]{CQUbadge.pdf}
		\caption{并排第一个图}
		\label{fig:parallel1}
	\end{minipage}\hfill
	\begin{minipage}{0.48\textwidth}
		\centering
		\includegraphics[height=5cm]{CQUbadge.pdf}
		\caption{并排第二个图}
		\label{fig:parallel2}
	\end{minipage}
\end{figure}

测试用途：theequation值为：\theequation ，thefigure值为：\thefigure ，thetable值为：\thetable

\chapter{The conclusion}


%\include{contents/yourFreeChoise}


\backmatter %%% 后置部分(致谢、参考文献、附录等)


%% 致谢
% Copyright (c) 2014,2016 Casper Ti. Vector
% Public domain.

\chapter{致谢}
\pkuthssffaq % 中文测试文字。

% vim:ts=4:sw=4

%% 参考文献
% 顺序编码制：cqunumerical		
% 注意：至少需要引用一篇参考文献，否则下面两行会引起编译错误。
\bibliographystyle{cqunumerical}
\bibliography{ref/refs}


%% 附录(按ABC...分节，证明、推导、程序、个人简历等)
\appendix

% 个人简历
\chapter{The first appendix}
\section{a section}

\chapter{The second appendix}
\section{a section}
\subsection{a subsection}

%% 原创声明和授权说明书，可选：用扫描页替换
%\cquauthpage[contents/authscan.pdf]
\cquauthpage

\end{document}
